\section*{Chapitre \ref{ch:g10:sport-and-health} --- Activité physique et santé}

\begin{solution}{exo:g10:sport-and-health:1}
Un cœur plus gros, avec un volume d'éjection plus grand et une
fréquence de repos plus basse~; plus de capillaires et de \omterm{def:g10:cells-common-unit:organelle}{mitochondries}
dans les \omterm{def:g10:muscles-and-joints:muscle}{fibres musculaires}~; plus de globules rouges~; des réserves de
\omterm{def:g10:exercise-and-energy:glycogen}{glycogène} plus grandes (et une meilleure régulation de la glycémie et
de la pression artérielle).
\end{solution}

\begin{solution}{exo:g10:sport-and-health:2}
Le volume d'éjection a grandi, si bien que les mêmes \qty{5}{L/min}
sont délivrés en moins de battements~: $Q = f \times V_s$ avec un $V_s$
plus grand demande un $f$ plus petit.
\end{solution}

\begin{solution}{exo:g10:sport-and-health:3}
Environ une heure par jour d'activité modérée. Non~: la marche, le vélo
pour se rendre en cours, les escaliers et le ménage comptent tous ---
ce qui importe, c'est le mouvement qui élève la \omterm{def:g10:exercise-and-energy:expenditure}{dépense énergétique}.
\end{solution}

\begin{solution}{exo:g10:sport-and-health:4}
L'usage de substances ou de méthodes qui élèvent la performance en
passant outre les régulations du corps, au lieu de bâtir une
adaptation. L'EPO passe outre le contrôle du nombre de globules
rouges~; les stimulants passent outre la fatigue, l'avertissement qui
protège le cœur.
\end{solution}

\begin{solution}{exo:g10:sport-and-health:5}
Fatigue persistante, fréquence cardiaque de repos qui remonte,
performances qui baissent (mais aussi sommeil perturbé, infections
fréquentes, douleurs de \omterm{def:g10:muscles-and-joints:muscle}{tendon} ou d'os).
\end{solution}

\begin{solution}{exo:g10:sport-and-health:6}
Après 12 semaines~: $72 - 60 = 12$ battements pour quatre séances,
$72 - 64 = 8$ pour deux. L'adaptation est proportionnelle à la
sollicitation --- plus de séances, plus de changement --- et dans les
deux groupes elle ralentit à l'approche du nouvel état.
\end{solution}

\begin{solution}{exo:g10:sport-and-health:7}
$V_s = Q/f$~: de $5000/72 \approx \qty{69}{mL}$ à $5000/54 \approx
\qty{93}{mL}$, soit un facteur $72/54 \approx 1.33$.
\end{solution}

\begin{solution}{exo:g10:sport-and-health:8}
$1 - 0.55 = 45\%$ de moins. Le premier palier, du moins actif au peu
actif (1.00 à 0.80), apporte le plus grand gain.
\end{solution}

\begin{solution}{exo:g10:sport-and-health:9}
\qty{2.0}{kg}, soit $2.0/55 \approx 3.6\%$ de sa masse. Son volume
sanguin baisse, sa fréquence cardiaque monte pour tenir le débit, sa
température corporelle grimpe, sa performance baisse~; les crampes
deviennent probables.
\end{solution}

\begin{solution}{exo:g10:sport-and-health:10}
$10 \times 1.1^7 \approx \qty{19.5}{km}$, de sorte que les
\qty{20}{km} sont atteints la huitième semaine~: environ deux mois.
\end{solution}

\begin{solution}{exo:g10:sport-and-health:11}
Chaque fibre s'épaissit en ajoutant des myofibrilles, et le système
nerveux apprend à recruter davantage de fibres à la fois et de façon
mieux synchronisée~; les deux élèvent la force. Le nombre de fibres est
fixé à l'âge adulte.
\end{solution}

\begin{solution}{exo:g10:sport-and-health:12}
L'oxygène par litre augmente d'un facteur $56/45 \approx 1.24$, soit
environ 24~\%~; à \omterm{def:g10:heart-lungs-effort:output}{débit cardiaque} et extraction inchangés, la
$\dot V\!\mathrm{O_2}$max augmente au plus de la même fraction. Mais un
sang à 56~\% de \omterm{def:g10:cells-common-unit:cell}{cellules} est bien plus épais~; ralenti par un cœur
endormi, il peut coaguler dans un vaisseau coronaire ou cérébral.
\end{solution}

\begin{solution}{exo:g10:sport-and-health:13}
En altitude, le corps élève lui-même son EPO en réponse à un manque
d'oxygène réel, d'une quantité contrôlée, et la rabaisse quand le
manque cesse~: la régulation est à l'œuvre. L'EPO injectée impose un
nombre que la régulation ne choisirait jamais, sans manque qui le
justifie et sans frein.
\end{solution}

\begin{solution}{exo:g10:sport-and-health:14}
Le muscle au travail prélève le glucose sans insuline, et le muscle
entraîné y reste plus sensible~; la marche abaisse donc directement la
glycémie et réduit l'insuline nécessaire à chaque repas. Un médicament
abaisse le sucre mais laisse la cause en place --- un muscle inutilisé
et insensible --- et ne fait rien pour le cœur, les os ni la masse
corporelle, que la marche améliore aussi.
\end{solution}

\begin{solution}{exo:g10:sport-and-health:15}
L'adaptation se construit pendant le repos entre des sollicitations
répétées et augmentées peu à peu~: un mois de charge intensive a plus
de chances de produire du surentraînement et une blessure que de la
forme, puisque le \omterm{def:g10:muscles-and-joints:muscle}{tendon} et l'os s'adaptent en mois et non en semaines.
Et les gains se démontent en quelques semaines d'arrêt~: rien n'est
«~acquis pour l'année~» par une seule salve. Une activité modérée et
régulière tout au long de l'année fait ce que le mois ne peut pas
faire.
\end{solution}

\begin{solution}{pb:g10:sport-and-health:1}
\textbf{1.} $4800/72 \approx \qty{67}{mL}$ avant~; $4800/56 \approx
\qty{86}{mL}$ après.

\textbf{2.} $42 \times 62 \approx \qty{2.6}{L/min}$~; $50 \times 62 =
\qty{3.1}{L/min}$~: un gain de 19~\%.

\textbf{3.} Débit maximal $\dot V\!\mathrm{O_2}/0.150$~: de \num{17.4}
à \qty{20.7}{L/min}~; volume d'éjection maximal à 200 battements~: de
\num{87} à \qty{103}{mL}.

\textbf{4.} Le cœur plus gros et plus fort explique le volume
d'éjection~; davantage de capillaires et de \omterm{def:g10:cells-common-unit:organelle}{mitochondries} dans les
fibres expliquerait un gain d'extraction.

\textbf{5.} Un facteur 4 en 26 semaines fait en moyenne $4^{1/26}
\approx 1.055$, soit environ 5.5~\% par semaine --- dans la règle des 10~\%
($1.1^{26}
\approx 12$ en aurait autorisé bien davantage).

\textbf{6.} Le surentraînement~: fréquence cardiaque de repos qui
remonte, sommeil perturbé, infection, performances qui baissent,
douleur dans un tissu à réparation lente.

\textbf{7.} Une lésion de fatigue de l'os du tibia (fracture de fatigue
ou inflammation de sa surface)~: l'os s'adapte le plus lentement de
tous les tissus, et la charge a été doublée en quinze jours.

\textbf{8.} La récupération est incomplète d'une séance à l'autre~; le
corps est dans un état de tension permanente, et l'adaptation qui avait
abaissé la fréquence est en train de se défaire.

\textbf{9.} Une semaine de quasi-repos et de sommeil (l'adaptation se
construit pendant le repos)~; puis deux semaines à la moitié de la
distance précédente, sans travail de vitesse, en remontant d'un dixième
par semaine~; le tibia vu par un médecin et mis au repos jusqu'à
disparition de la douleur (écouter la douleur)~; boire et manger pour
la réparation.

\textbf{10.} Un stimulant~: il masque la fatigue, l'avertissement qui
protège le cœur de l'épuisement et de la surchauffe.

\textbf{11.} Sam $32 \times 80 \approx \qty{2.6}{L/min}$~; Alex $48
\times 64 \approx \qty{3.1}{L/min}$. Celle d'Alex est la plus grande
même en valeur absolue~; et la vie de tous les jours --- escaliers,
marche, se porter soi-même --- déplace la masse du corps, si bien que
le chiffre par kilogramme est celui qui décide de ce que l'on ressent.

\textbf{12.} Sam $12/32 \approx 38\%$ de son maximum, Alex $12/48 =
25\%$. Sam, qui travaille au-dessus du tiers de son plafond, arrive
essoufflé.

\textbf{13.} Un muscle inutilisé prélève peu de glucose dans le sang et
perd sa sensibilité à l'insuline~; les mêmes repas demandent plus
d'insuline chaque année et la glycémie commence à monter.

\textbf{14.} De 1.00 à 0.70~: environ 30~\% de risque en moins.

\textbf{15.} L'\omterm{def:g10:sport-and-health:training}{activité physique} est tout mouvement qui élève la
\omterm{def:g10:exercise-and-energy:expenditure}{dépense énergétique}~: marcher pour aller travailler, les escaliers au
lieu de l'ascenseur, le ménage --- une demi-heure par jour n'exige ni
sport ni matériel.

\textbf{16.} À 1~\% par an, une marge de 10~\% vaut dix ans de déclin.

\textbf{17.} Le sommet de Sam était plus bas et il perd de l'os depuis
vingt-cinq ans~; son poignet a franchi le seuil auquel une chute le
casse, alors que l'os d'Alex, bâti sous la charge à l'adolescence et
entretenu depuis, garde dix ans d'avance.

\textbf{18.} Une adaptation est une structure entretenue par l'usage,
non un acquis définitif~: elle est démontée quand la sollicitation
cesse et rebâtie quand elle reprend.

\textbf{19.} Maladies du cœur et des artères, diabète de type 2,
ostéoporose (mais aussi certains cancers et la dépression). Le
diabète~: un muscle inactif perd sa sensibilité à l'insuline, la
glycémie monte et le pancréas est surmené.

\textbf{20.} Fréquence cardiaque de repos 54 contre 78~;
$\dot V\!\mathrm{O_2}$max \num{48} contre \qty{32}{mL/min} par
kilogramme~; environ trois heures de course par semaine, tenues pendant
quatorze ans.
\end{solution}
